L'operateur de diffusion quasilineaire a ete developpe par Kennel
et Engelmann . Un traitement relativiste effectue par Lerche propose
un operateur qui peut etre reecrit sous la forme

\[
Q(f_{0})=\lim_{V\rightarrow\infty}\frac{\pi}{V}\int\frac{d^{3}k}{(2\pi)^{3}}\sum_{n=-\infty}^{+\infty}\mathcal{L}\left[p_{\bot}^{2}D_{\mathbf{k}}^{(n)}\delta(\omega-k_{\Vert}v_{\Vert}-n\Omega)\mathcal{L}(f_{0})\right]\]
 ou $\mathcal{L}$ est un operateur differentiel dans l'espace des
moments defini par \[
\mathcal{L}(\mathcal{A})=\frac{1}{p}\frac{\partial\mathcal{A}}{\partial p}+\frac{1}{p^{2}\xi}\left(\frac{k_{\Vert}v_{\Vert}}{\omega}-\xi^{2}\right)\frac{\partial\mathcal{A}}{\partial\xi}\]
 qui peut etre sous la forme de la divergence d'un flux \[
\mathcal{L}(\mathcal{A})=\frac{1}{p^{2}}\left[\frac{\partial}{\partial p}\left(p\mathcal{A}\right)+\frac{\partial}{\partial\xi}\left(\frac{1}{\xi}\left[\frac{k_{\Vert}v_{\Vert}}{\omega}-\xi^{2}\right]\mathcal{A}\right)\right]\]
 $D_{\mathbf{k}}$ etant le coefficient de diffusion quasilineaire,
qui prend en compte l'intensite et la polarisation du champ electrique
de l'onde RF. Il est defini par \[
D_{\mathbf{k}}^{(n)}=q_{e}^{2}\left|E_{\mathbf{k},+}e^{-i\alpha}J_{n-1}\left(\frac{k_{\bot}v_{\bot}}{\Omega}\right)+E_{\mathbf{k},-}e^{+i\alpha}J_{n+1}\left(\frac{k_{\bot}v_{\bot}}{\Omega}\right)+\frac{p_{\Vert}}{p_{\bot}}E_{\mathbf{k},\Vert}J_{n}\left(\frac{k_{\bot}v_{\bot}}{\Omega}\right)\right|^{2}\]
 avec les notations suivantes: le champ electrique est decompose en
ses composantes de Fourier \[
\overrightarrow{E}(\overrightarrow{x},t)=\iiint\frac{d^{3}k}{(2\pi)^{3}}e^{i\mathbf{k}\cdot\mathbf{x}}\overrightarrow{E}_{\mathbf{k}}\left(\mathbf{k},t\right)\]
 qui sont exprimees en fonction des coordonnees de champ tournant
\[
\left\{ \begin{array}{l}
E_{\mathbf{k},\pm}=\frac{1}{2}\left(E_{\mathbf{k},x}\pm iE_{\mathbf{k},y}\right)\\
E_{\mathbf{k},\Vert}=E_{\mathbf{k},z}\end{array}\right.\]
 l'angle $\alpha$ decrivant la projection de la composante transverse
du vecteur d'onde dans le plan $(x,y)$: \begin{eqnarray*}
k_{x} & = & k_{\bot}\cos\alpha\\
k_{y} & = & k_{\bot}\sin\alpha\\
k_{z} & = & k_{\Vert}\end{eqnarray*}


En permutant les operateurs differentiels dans l'espace des moments
et l'operateur integral dans l'espace des vecteurs d'onde $k$, $Q(f_{0})$
devient \begin{eqnarray*}
Q(f_{0}) & = & \sum_{n=-\infty}^{+\infty}\frac{1}{p^{2}}\frac{\partial}{\partial p}\left(p\lim_{V\rightarrow\infty}\frac{\pi}{V}\int\frac{d^{3}k}{(2\pi)^{3}}p_{\bot}^{2}D_{\mathbf{k}}^{(n)}\delta(\omega-k_{\Vert}v_{\Vert}-n\Omega)\mathcal{L}(f_{0})\right)\\
 &  & +\frac{1}{p^{2}}\frac{\partial}{\partial\xi}\left(\frac{1}{\xi}\lim_{V\rightarrow\infty}\frac{\pi}{V}\int\frac{d^{3}k}{(2\pi)^{3}}\left[\frac{k_{\Vert}v_{\Vert}}{\omega}-\xi^{2}\right]p_{\bot}^{2}D_{\mathbf{k}}^{(n)}\delta(\omega-k_{\Vert}v_{\Vert}-n\Omega)\mathcal{L}(f_{0})\right)\end{eqnarray*}


En exprimant la condition de resonance \[
k_{\Vert}v_{\Vert}=\omega-n\Omega\]
 on obtient un operateur quasilineaire sous la forme de la divergence
d'un flux \[
Q(f_{0})=-\nabla\cdot\mathbf{S}^{\text{RF}}=-\frac{1}{p^{2}}\frac{\partial}{\partial p}\left(p^{2}S_{p}^{\text{RF}}\right)+\frac{1}{p}\frac{\partial}{\partial\xi}\left(\sqrt{1-\xi^{2}}S_{\xi}^{\text{RF}}\right)\]
 pour lequel les contributions de chaque rayon sont separees \begin{eqnarray*}
S_{p}^{\text{RF}} & = & \sum_{n=-\infty}^{+\infty}\sum_{b}-p(1-\xi^{2})D_{b,n}^{\text{RF}}(p,\xi)\mathcal{L}(f_{0})\\
S_{\xi}^{\text{RF}} & = & \sum_{n=-\infty}^{+\infty}\sum_{b}p\frac{\sqrt{1-\xi^{2}}}{\xi}\left[1-\xi^{2}-\frac{n\Omega}{\omega_{b}}\right]D_{b,n}^{\text{RF}}(p,\xi)\mathcal{L}(f_{0})\end{eqnarray*}
 et \[
D_{b,n}^{\text{RF}}(p,\xi)=\lim_{V\rightarrow\infty}\frac{\pi}{V}\int\frac{d^{3}k}{(2\pi)^{3}}D_{\mathbf{k}}^{b(n)}\delta(\omega_{b}-k_{\Vert}v_{\Vert}-n\Omega)\]


De plus, les flux quasilineaires sont purement diffusifs et ils peuvent
etre exprimes sous la forme \[
\mathbf{S}^{\text{RF}}=-\overline{\overline{D}}_{\mathbf{p}}^{\text{RF}}\cdot\nabla_{\mathbf{p}}f_{0}\]
 avec \[
\left\{ \begin{array}{l}
D_{pp}^{\text{RF}}=\sum_{n=-\infty}^{+\infty}\sum_{b}(1-\xi^{2})D_{b,n}^{\text{RF}}(p,\xi)\\
D_{p\xi}^{\text{RF}}=\sum_{n=-\infty}^{+\infty}\sum_{b}-\frac{\sqrt{1-\xi^{2}}}{\xi}\left[1-\xi^{2}-\frac{n\Omega}{\omega_{b}}\right]D_{b,n}^{\text{RF}}(p,\xi)\\
D_{\xi p}^{\text{RF}}=\sum_{n=-\infty}^{+\infty}\sum_{b}-\frac{\sqrt{1-\xi^{2}}}{\xi}\left[1-\xi^{2}-\frac{n\Omega}{\omega_{b}}\right]D_{b,n}^{\text{RF}}(p,\xi)\\
D_{\xi\xi}^{\text{RF}}=\sum_{n=-\infty}^{+\infty}\sum_{b}\frac{1}{\xi^{2}}\left[1-\xi^{2}-\frac{n\Omega}{\omega_{b}}\right]^{2}D_{b,n}^{\text{RF}}(p,\xi)\end{array}\right.\]


A partir de ces expressions il est facile de determiner la moyenne
des coefficients sur la periode de rebond, \[
\left\{ \begin{array}{l}
D_{pp}^{\text{RF}(0)}=\sum_{n=-\infty}^{+\infty}\sum_{b}(1-\xi_{0}^{2})\left\{ \Psi D_{(n)}^{\text{RF}}(p,\xi)\right\} \\
D_{p\xi}^{\text{RF}(0)}=\sum_{n=-\infty}^{+\infty}\sum_{b}-\frac{\sqrt{1-\xi_{0}^{2}}}{\xi_{0}}\left[1-\xi_{0}^{2}-\frac{n\Omega_{0}}{\omega_{b}}\right]\left\{ \Psi D_{(n)}^{\text{RF}}(p,\xi)\right\} \\
D_{\xi p}^{\text{RF}(0)}=\sum_{n=-\infty}^{+\infty}\sum_{b}-\frac{\sqrt{1-\xi_{0}^{2}}}{\xi_{0}}\left[1-\xi_{0}^{2}-\frac{n\Omega_{0}}{\omega_{b}}\right]\left\{ \Psi D_{(n)}^{\text{RF}}(p,\xi)\right\} \\
D_{\xi\xi}^{\text{RF}(0)}=\sum_{n=-\infty}^{+\infty}\sum_{b}\frac{1}{\xi_{0}^{2}}\left[1-\xi_{0}^{2}-\frac{n\Omega_{0}}{\omega_{b}}\right]^{2}\left\{ \Psi D_{(n)}^{\text{RF}}(p,\xi)\right\} \end{array}\right.\]
 ou l'on introduit la frequence cyclotronique au demi-plan $\theta=0$:
\[
\Omega_{0}=\frac{q_{e}B_{0}}{\gamma m_{e}}=\frac{\Omega}{\Psi}\]


On peut finalement reecrire, \[
\left\{ Q(f_{0})\right\} =\left\{ -\nabla\cdot\mathbf{S}^{\text{RF}}\right\} =-\frac{1}{p^{2}}\frac{\partial}{\partial p}\left(p^{2}S_{p}^{\text{RF}(0)}\right)+\frac{1}{\lambda\left(\xi_{0}\right)p}\frac{\partial}{\partial\xi}\left(\sqrt{1-\xi^{2}}\lambda\left(\xi_{0}\right)S_{\xi}^{\text{RF}(0)}\right)\]
 avec \[
\mathbf{S}^{\text{RF}(0)}=-\overline{\overline{D}}_{\mathbf{p}}^{\text{RF}(0)}\cdot\nabla_{p,\xi_{0}}f_{0}^{(0)}\]
 et \[
\left\{ \begin{array}{l}
D_{pp}^{\text{RF}(0)}=\sum_{n=-\infty}^{+\infty}\sum_{b}(1-\xi_{0}^{2})D_{b,n}^{\text{RF}(0)}(p,\xi_{0})\\
D_{p\xi}^{\text{RF}(0)}=\sum_{n=-\infty}^{+\infty}\sum_{b}-\frac{\sqrt{1-\xi_{0}^{2}}}{\xi_{0}}\left[1-\xi_{0}^{2}-\frac{n\Omega_{0}}{\omega}\right]D_{b,n}^{\text{RF}(0)}(p,\xi_{0})\\
D_{\xi p}^{\text{RF}(0)}=\sum_{n=-\infty}^{+\infty}\sum_{b}-\frac{\sqrt{1-\xi_{0}^{2}}}{\xi_{0}}\left[1-\xi_{0}^{2}-\frac{n\Omega_{0}}{\omega_{b}}\right]D_{b,n}^{\text{RF}(0)}(p,\xi_{0})\\
D_{\xi\xi}^{\text{RF}(0)}=\sum_{n=-\infty}^{+\infty}\sum_{b}\frac{1}{\xi_{0}^{2}}\left[1-\xi_{0}^{2}-\frac{n\Omega_{0}}{\omega_{b}}\right]^{2}D_{b,n}^{\text{RF}(0)}(p,\xi_{0})\end{array}\right.\]
 ou par definition \[
D_{b,n}^{\text{RF}(0)}(p,\xi_{0})=\left\{ \Psi D_{b,n}^{\text{RF}}(p,\xi)\right\} \]


La discretisation des elements du tenseur quasi-lineaire est alors
immediate, et les coefficients qui interviennent dans la matrice $\overline{\overline{M}}^{\left(0\right)}$sont
\[
\overline{\overline{D}}_{\mathbf{p}}^{\text{RF}(0)}\rightarrow\left\{ \begin{array}{l}
D_{pp,i+1,j+1/2}^{\text{RF}\left(0\right)}=\sum_{n=-\infty}^{+\infty}\sum_{b}(1-\xi_{0j+1/2}^{2})D_{b,n,i+1,j+1/2}^{\text{ RF}(0)}\\
D_{pp,i,j+1/2}^{\text{RF}\left(0\right)}=\sum_{n=-\infty}^{+\infty}\sum_{b}(1-\xi_{0j+1/2}^{2})D_{b,n,i,j+1/2}^{\text{RF}(0)}\\
D_{p\xi,i+1,j+1/2}^{\text{RF}\left(0\right)}=\sum_{n=-\infty}^{+\infty}\sum_{b}-\sqrt{1-\xi_{0j+1/2}^{2}}/\xi_{0j+1/2}\left[1-\xi_{0j+1/2}^{2}-n\Omega_{0}/\omega_{b}\right]D_{b,n,i+1,j+1/2}^{\text{RF}(0)}\\
D_{\xi p,i+1/2,j+1}^{\text{RF}\left(0\right)}=\sum_{n=-\infty}^{+\infty}\sum_{b}-\sqrt{1-\xi_{0j+1}^{2}}/\xi_{0j+1}\left[1-\xi_{0j+1}^{2}-n\Omega_{0}/\omega_{b}\right]D_{b,n,i+1/2,j+1}^{\text{RF}(0)}\\
D_{p\xi,i+1/2,j+1/2}^{\text{RF}\left(0\right)}=\sum_{n=-\infty}^{+\infty}\sum_{b}-\sqrt{1-\xi_{0j+1/2}^{2}}/\xi_{0j+1/2}\left[1-\xi_{0j+1/2}^{2}-n\Omega_{0}/\omega_{b}\right]D_{b,n,i+1/2,j+1/2}^{\text{RF}(0)}\\
D_{\xi p,i+1/2,j+1/2}^{\text{RF}\left(0\right)}=\sum_{n=-\infty}^{+\infty}\sum_{b}-\sqrt{1-\xi_{0j+1/2}^{2}}/\xi_{0j+1/2}\left[1-\xi_{0j+1/2}^{2}-n\Omega_{0}/\omega_{b}\right]D_{b,n,i+1/2,j+1/2}^{\text{RF}(0)}\\
D_{p\xi,i,j+1/2}^{\text{RF}\left(0\right)}=\sum_{n=-\infty}^{+\infty}\sum_{b}-\sqrt{1-\xi_{0j+1/2}^{2}}/\xi_{0j+1/2}\left[1-\xi_{0j+1/2}^{2}-n\Omega_{0}/\omega_{b}\right]D_{b,n,i,j+1/2}^{\text{RF}(0)}\\
D_{\xi p,i+1/2,j}^{\text{RF}\left(0\right)}=\sum_{n=-\infty}^{+\infty}\sum_{b}-\sqrt{1-\xi_{0j}^{2}}/\xi_{0j}\left[1-\xi_{0j}^{2}-n\Omega_{0}/\omega_{b}\right]D_{b,n,i+1/2,j}^{\text{RF}(0)}\\
D_{\xi\xi,i+1/2,j+1}^{\text{RF}\left(0\right)}=\sum_{n=-\infty}^{+\infty}\sum_{b}1/\xi_{0j+1}^{2}\left[1-\xi_{0j+1}^{2}-n\Omega_{0}/\omega_{b}\right]^{2}D_{b,n,i+1/2,j+1}^{\text{RF}(0)}\\
D_{\xi\xi,i+1/2,j}^{\text{RF}\left(0\right)}=\sum_{n=-\infty}^{+\infty}\sum_{b}1/\xi_{0j}^{2}\left[1-\xi_{0j}^{2}-n\Omega_{0}/\omega_{b}\right]^{2}D_{b,n,i+1/2,j}^{\text{RF}(0)}\end{array}\right.\]
 et \[
\overline{\overline{F}}_{p}^{\text{RF}(0)}\rightarrow\left\{ \begin{array}{l}
F_{p,i+1,j+1/2}^{\text{RF}(0)}=0\\
F_{p,i,j+1/2}^{\left(0\right)}=0\\
F_{\xi,i+1/2,j+1}^{\left(0\right)}=0\\
F_{\xi,i+1/2,j}^{\left(0\right)}=0\end{array}\right.\]


Sous cette forme, le calcul garde une tres grande generalite, quelque
soit le type d'onde.


